\bibitem{vapnik1971uniform} Vapnik, Vladimir N, Chervonenkis, Alexey Y. On the uniform convergence of relative frequencies of events to their probabilities. Theory of Probability and its Applications. V.~16, N~2. P.~264--280. 1971.
\bibitem{valiant1984theory} Valiant, Leslie G. A theory of the learnable. Communications of the ACM. V.~27, N~11. P.~1134--1142. 1984.
\bibitem{vapnik1998statistical} Vapnik, Vladimir N. Statistical Learning Theory. Wiley. 1998.
\bibitem{vorontsov2004combinatorial} Воронцов, Константин В. Комбинаторный подход к оценке качества обучаемых алгоритмов. Математические вопросы кибернетики. V.~13. P.~5--51. 2004.
\bibitem{bartlett2002rademacher} Bartlett, Peter L, Mendelson, Shahar. Rademacher and Gaussian complexities: Risk bounds and structural results. Journal of Machine Learning Research. V.~3. P.~463--482. 2002.
\bibitem{koltchinskii2001rademacher} Koltchinskii, Vladimir. Rademacher penalties and structural risk minimization. IEEE Transactions on Information Theory. V.~47, N~5. P.~1902--1914. 2001.
\bibitem{mcallester1999pac} McAllester, David A. Some PAC-Bayesian theorems. Machine Learning. V.~37, N~3. P.~355--363. 1999.
\bibitem{dziugaite2017computing} Dziugaite, Gintare Karolina, Roy, Daniel M. Computing nonvacuous generalization bounds for deep (stochastic) neural networks with many more parameters than training data. Uncertainty in Artificial Intelligence. 2017.
\bibitem{xu2017information} Xu, Aolin, Raginsky, Maxim. Information-theoretic analysis of generalization capability of learning algorithms. Advances in Neural Information Processing Systems. V.~30. 2017.
\bibitem{bousquet2002stability} Bousquet, Olivier, Elisseeff, Andre. Stability and generalization. Journal of Machine Learning Research. V.~2. P.~499--526. 2002.
\bibitem{hardt2016train} Hardt, Moritz, Recht, Benjamin, Singer, Yoram. Train faster, generalize better: stability of stochastic gradient descent. Proceedings of the 33rd International Conference on Machine Learning. P.~1225--1234. 2016.
\bibitem{li2018visualizing} Li, Hao, Xu, Zheng, Taylor, Gavin, Studer, Christoph, Goldstein, Tom. Visualizing the Loss Landscape of Neural Nets. Advances in Neural Information Processing Systems. V.~31. 2018.
\bibitem{hochreiter1997flat} Hochreiter, Sepp, Schmidhuber, Jurgen. Flat minima. Neural Computation. V.~9, N~1. P.~1--42. 1997.
\bibitem{sagun2017eigenvalues} Sagun, Levent, Bottou, Leon, LeCun, Yann. Eigenvalues of the Hessian in Deep Learning: Singularity and Beyond. arXiv preprint arXiv:1611.07476. 2017.
\bibitem{kiselev2024unraveling} Kiselev, Nikita, Grabovoy, Andrey. Unraveling the Hessian: A Key to Smooth Convergence in Loss Function Landscapes. Doklady Mathematics. V.~110, N~S1. P.~S42. 2025.
\bibitem{belkin2019reconciling} Belkin, Mikhail, Hsu, Daniel, Ma, Siyuan, Mandal, Soumik. Reconciling modern machine-learning practice and the classical bias--variance trade-off. Proceedings of the National Academy of Sciences. V.~116, N~32. P.~15849--15854. 2019.
\bibitem{jacot2018neural} Jacot, Arthur, Gabriel, Franck, Hongler, Clement. Neural tangent kernel: Convergence and generalization in neural networks. Advances in Neural Information Processing Systems. V.~31. 2018.
\bibitem{kaplan2020scaling} Kaplan, Jared, McCandlish, Sam, Henighan, Tom, Brown, Tom B, Chess, Benjamin, Child, Rewon, Gray, Scott, Radford, Alec, Wu, Jeffrey, Amodei, Dario. Scaling laws for neural language models. arXiv preprint arXiv:2001.08361. 2020.
\bibitem{hoffman2022training} Hoffmann, Jordan, Borgeaud, Sebastian, Mensch, Arthur, Buchatskaya, Elena, Cai, Trevor, Rutherford, Eliza, de Las Casas, Diego, Hendricks, Lisa Anne, Welbl, Johannes, Clark, Aidan, Hennigan, Tom, Noland, Eric, Millican, Katie, van den Driessche, George, Damoc, Bogdan, Guy, Aurelia, Osindero, Simon, Simonyan, Karen, Elsen, Erich, Vinyals, Oriol, Rae, Jack W., Sifre, Laurent. Training compute-optimal large language models. Proceedings of the 36th International Conference on Neural Information Processing Systems. Curran Associates Inc. 2022.
\bibitem{vaswani2017attention} Vaswani, Ashish, Shazeer, Noam, Parmar, Niki, Uszkoreit, Jakob, Jones, Llion, Gomez, Aidan N, Kaiser, Lukasz, Polosukhin, Illia. Attention is all you need. Advances in Neural Information Processing Systems. V.~30. 2017.
\bibitem{ho2020denoising} Ho, Jonathan, Jain, Ajay, Abbeel, Pieter. Denoising diffusion probabilistic models. Advances in Neural Information Processing Systems. V.~33. P.~6840--6851. 2020.
\bibitem{hellstrom2025perspectives} Hellstrom, Fredrik, Durisi, Giuseppe, Guedj, Benjamin, Raginsky, Maxim. Generalization Bounds: Perspectives from Information Theory and PAC-Bayes. Foundations and Trends in Machine Learning. V.~18, N~4. P.~433--785. 2025.
\bibitem{petrov2025curvaturegap} Petrov, Egor, Kiselev, Nikita, Meshkov, Vladislav, Grabovoy, Andrey. Closing the Curvature Gap: Full Transformer Hessians and Their Implications for Scaling Laws. arXiv preprint arXiv:2504.12516. 2025.
\bibitem{transformerScalingLowDim2024} Understanding Scaling Laws with Statistical and Approximation Theory for Transformer Neural Networks on Intrinsically Low-dimensional Data. arXiv. 2024.
\bibitem{transformerScalingStatApprox2022} Scaling Laws for Transformers on Low-Dimensional Data: A Statistical and Approximation Theory Perspective. arXiv. 2022.
\bibitem{moeScalingLaws2026} Generalization and Scaling Laws for Mixture-of-Experts Transformers. arXiv. 2026.
\bibitem{featureDynamicsScaling2025} Understanding Scaling Laws in Deep Neural Networks via Feature Learning Dynamics. arXiv. 2025.
\bibitem{diffusionTransformerScaling2024} Scaling Laws For Diffusion Transformers. arXiv. 2024.
\bibitem{beyondChinchillaInference2023} Beyond Chinchilla-Optimal: Accounting for Inference in Language Model Scaling Laws. arXiv. 2023.
\bibitem{zhang2021rethinking} Zhang, Chiyuan, Bengio, Samy, Hardt, Moritz, Recht, Benjamin, Vinyals, Oriol. Understanding deep learning (still) requires rethinking generalization. Communications of the ACM. V.~64, N~3. P.~107--115. 2021.
\bibitem{truong2022rademacher} Lan V. Truong. On Rademacher Complexity-based Generalization Bounds for Deep Learning. arXiv. 2022.
\bibitem{algorithmicRademacher2023} Sarah Sachs; Tim van Erven; Liam Hodgkinson; Rajiv Khanna; Umut Simsekli. Generalization Guarantees via Algorithm-dependent Rademacher Complexity. arXiv. 2023.
\bibitem{zhou2018pacbayescompression} Wenda Zhou; Victor Veitch; Morgane Austern; Ryan P. Adams; Peter Orbanz. Non-Vacuous Generalization Bounds at the ImageNet Scale: A PAC-Bayesian Compression Approach. arXiv (Cornell University). 2018.
\bibitem{pacbayesinterpolating2023} Liam Hodgkinson; Chris van der Heide; Robert Salomone; Fred Roosta; Michael W. Mahoney. A PAC-Bayesian Perspective on the Interpolating Information Criterion. arXiv. 2023.
\bibitem{arora2019finegrained} Arora, Sanjeev, Du, Simon, Hu, Wei, Li, Zhiyuan, Wang, Ruosong. Fine-grained analysis of optimization and generalization for overparameterized two-layer neural networks. International Conference on Machine Learning. P.~322--332. 2019.
\bibitem{adlam2020ntkhighdim} Ben Adlam; Jeffrey Pennington. The Neural Tangent Kernel in High Dimensions: Triple Descent and a Multi-Scale Theory of Generalization. arXiv (Cornell University). 2020.
\bibitem{keskar2017largebatch} Keskar, Nitish Shirish, Mudigere, Dheevatsa, Nocedal, Jorge, Smelyanskiy, Mikhail, Tang, Ping. On Large-Batch Training for Deep Learning: Generalization Gap and Sharp Minima. arXiv preprint arXiv:1609.04836. 2017.
\bibitem{foret2021sam} Foret, Pierre, Kleiner, Ariel, Mobahi, Hossein, Neyshabur, Behnam. Sharpness-Aware Minimization for Efficiently Improving Generalization. International Conference on Machine Learning. PMLR. V.~119. P.~7295--7305. 2020.
\bibitem{wu2020sgdflat} Zeke Xie; Issei Sato; Masashi Sugiyama. A Diffusion Theory For Deep Learning Dynamics: Stochastic Gradient Descent Exponentially Favors Flat Minima. arXiv. 2020.
\bibitem{wu2017landscapes} Lei Wu; Zhanxing Zhu; E Weinan. Towards Understanding Generalization of Deep Learning: Perspective of Loss Landscapes. arXiv (Cornell University). 2017.
\bibitem{song2019randomfeatures} Mei Song; Andrea Montanari. The generalization error of random features regression: Precise asymptotics and double descent curve. arXiv (Cornell University). 2019.
\bibitem{valleperez2020bounds} Guillermo Valle-Perez; Ard A. Louis. Generalization bounds for deep learning. arXiv. 2020.
\bibitem{hu2021complexitysurvey} Xia Hu; Lingyang Chu; Jian Pei; Weiqing Liu; Jiang Bian. Model Complexity of Deep Learning: A Survey. arXiv. 2021.
\bibitem{recent2020theory} Fengxiang He; Dacheng Tao. Recent advances in deep learning theory. arXiv. 2020.
\bibitem{advani2020highdim} Madhu Advani; Andrew Saxe; Haim Sompolinsky. High-dimensional dynamics of generalization error in neural networks. Neural Networks. 2020.
\bibitem{vapnik1999vccontrol} Vladimir Cherkassky; Xuhui Shao; Filip Mulier; Vladimir Vapnik. Model complexity control for regression using VC generalization bounds. IEEE Transactions on Neural Networks. 1999.
\bibitem{bartlett1996vcreal} Arne Hole. Vapnik-Chervonenkis Generalization Bounds for Real Valued Neural Networks. Neural Computation. 1996.
\bibitem{belkin2021overview} Yehuda Dar; Vidya Muthukumar; Richard G. Baraniuk. A Farewell to the Bias-Variance Tradeoff? An Overview of the Theory of Overparameterized Machine Learning. arXiv (Cornell University). 2021.
\bibitem{nakkiran2021doubledescent} Nakkiran, Preetum, Kaplun, Gal, Bansal, Yamini, Yang, Tristan, Barak, Boaz, Sutskever, Ilya. Deep double descent: Where bigger models and more data hurt. Journal of Statistical Mechanics: Theory and Experiment. V.~2021, N~12. P.~124003. 2021.
\bibitem{shcherbina2024vcdd} Eng Hock Lee; Vladimir Cherkassky. Understanding Double Descent Using VC-Theoretical Framework. IEEE Transactions on Neural Networks and Learning Systems. 2024.
\bibitem{pacbayes2025survey} Ghojogh, Benyamin, Ghodsi, Ali. PAC Learnability and Information Bottleneck in Deep Learning: Tutorial and Survey. arXiv preprint. 2024.
\bibitem{implicit2026synthesis} Zeran Johannsen. Implicit Regularization and Generalization in Overparameterized Neural Networks. arXiv. 2026.
\bibitem{uniform2021ntk} Sattar Vakili; Michael Bromberg; Jezabel Garcia; Da-shan Shiu; Alberto Bernacchia. Uniform Generalization Bounds for Overparameterized Neural Networks. arXiv. 2021.
\bibitem{leaveoneout2022} Gregor Bachmann; Thomas Hofmann; Aurelien Lucchi. Generalization Through The Lens Of Leave-One-Out Error. arXiv. 2022.
\bibitem{featurevskernel2020} Taiji Suzuki; Shunta Akiyama. Benefit of deep learning with non-convex noisy gradient descent: Provable excess risk bound and superiority to kernel methods. arXiv. 2020.
\bibitem{lottery2022pacbayes} Keitaro Sakamoto; Issei Sato. Analyzing Lottery Ticket Hypothesis from PAC-Bayesian Theory Perspective. arXiv (Cornell University). 2022.
\bibitem{normalizedflat2019} Yusuke Tsuzuku; Issei Sato; Masashi Sugiyama. Normalized Flat Minima: Exploring Scale Invariant Definition of Flat Minima for Neural Networks using PAC-Bayesian Analysis. arXiv. 2019.
\bibitem{hessian2019sgd} Mingwei Wei; David J. Schwab. How noise affects the Hessian spectrum in overparameterized neural networks. arXiv (Cornell University). 2019.
\bibitem{sam2023eos} Philip M. Long; Peter L. Bartlett. Sharpness-Aware Minimization and the Edge of Stability. arXiv. 2023.
\bibitem{hessianpowerlaw2022} Zeke Xie; Qian-Yuan Tang; Yunfeng Cai; Mingming Sun; Ping Li. On the Power-Law Hessian Spectrums in Deep Learning. arXiv (Cornell University). 2022.
\bibitem{sgdstability2023} Lei Wu; Weijie J. Su. The Implicit Regularization of Dynamical Stability in Stochastic Gradient Descent. arXiv. 2023.
\bibitem{landscape2017algorithms} Pan Zhou; Jiashi Feng. The Landscape of Deep Learning Algorithms. arXiv (Cornell University). 2017.
\bibitem{representation2020complexity} Parth Natekar; Manik Sharma. Representation Based Complexity Measures for Predicting Generalization in Deep Learning. arXiv. 2020.
\bibitem{simplicitybias2018} Guillermo Valle-Perez; Chico Q. Camargo; Ard A. Louis. Deep learning generalizes because the parameter-function map is biased towards simple functions. arXiv (Cornell University). 2018.
\bibitem{sgd2020stability} Yunwen Lei; Yiming Ying. Fine-Grained Analysis of Stability and Generalization for Stochastic Gradient Descent. arXiv (Cornell University). 2020.
\bibitem{pacbayes2019sgd} Ben London. A PAC-Bayesian Analysis of Randomized Learning with Application to Stochastic Gradient Descent. Neural Information Processing Systems. 2017.
\bibitem{lowpass2022flat} Devansh Bisla; Jing Wang; Anna Choromanska. Low-Pass Filtering SGD for Recovering Flat Optima in the Deep Learning Optimization Landscape. arXiv (Cornell University). 2022.
\bibitem{fractal2019sgd} Alexander Camuto; George Deligiannidis; Murat A. Erdogdu; Mert Gurbuzbalaban; Umut Simsekli; Lingjiong Zhu. Fractal Structure and Generalization Properties of Stochastic Optimization Algorithms. arXiv. 2021.
\bibitem{meta2024survey} Jing Wang; Anna Choromanska. A Survey of Optimization Methods for Training DL Models: Theoretical Perspective on Convergence and Generalization. arXiv. 2025.
\bibitem{explaining2021thesis} Vaishnavh Nagarajan. Explaining generalization in deep learning: progress and fundamental limits. arXiv. 2021.
\bibitem{kawaguchi2017generalization} Kenji Kawaguchi; Leslie Pack Kaelbling; Yoshua Bengio. Generalization in Deep Learning. arXiv. 2017.
\bibitem{mei2021eos} Sanjeev Arora; Zhiyuan Li; Abhishek Panigrahi. Understanding Gradient Descent on Edge of Stability in Deep Learning. arXiv. 2022.
\bibitem{poggio2020theoretical} Tomaso Poggio; Andrzej Banburski; Qianli Liao. Theoretical Issues in Deep Networks: Approximation, Optimization and Generalization. arXiv. 2019.
\bibitem{poggio2017nonoverfitting} Tomaso Poggio; Kenji Kawaguchi; Qianli Liao; Brando Miranda; Lorenzo Rosasco; Xavier Boix; Jack Hidary; Hrushikesh Mhaskar. Theory of Deep Learning III: explaining the non-overfitting puzzle. arXiv. 2017.
\bibitem{golowich2016lessons} Golowich, Noah, Rakhlin, Alexander, Shamir, Ohad. Size-Independent Sample Complexity of Neural Networks. Journal of Machine Learning Research. V.~19, N~1. P.~2978--3019. 2018.
\bibitem{seleznova2022ntklimits} Nikhil Vyas; Yamini Bansal; Preetum Nakkiran. Limitations of the NTK for Understanding Generalization in Deep Learning. arXiv. 2022.
